% TOP_PROBLEM: {"id": 3321, "title": "Drawing disconnected graphs on surfaces", "category_id": 3, "category": {"id": 3, "name": "graph_theory", "display_name": "Graph Theory", "description": "Problems involving graphs, networks, and their properties.", "slug": "graph-theory", "order_index": 3, "created_at": "2026-07-31T15:26:25.671Z"}, "status": "open", "problem_number": "OPG-322", "set_id": 12}
% FIRST_POSTED: 2026-10-01
% ENTRY_KIND: statement_only
% UnsolvedMath ID 3321; source catalogue code OPG-322
% !TeX program = lualatex
% Definition and sources only; this file is not an LLM solution attempt.
\documentclass[11pt,a4paper]{article}
\usepackage[margin=25mm]{geometry}
\usepackage{fontspec}
\setmainfont{lmroman10-regular.otf}[BoldFont=lmroman10-bold.otf,
 ItalicFont=lmroman10-italic.otf,BoldItalicFont=lmroman10-bolditalic.otf]
\usepackage{amsmath,amssymb,mathtools}
\usepackage{unicode-math}
\setmathfont{latinmodern-math.otf}
\AtBeginDocument{\let\setminus\smallsetminus}
\usepackage{xurl,hyperref}
\hypersetup{colorlinks=true,urlcolor=blue}
\setcounter{secnumdepth}{0}
\setlength{\parindent}{0pt}
\setlength{\parskip}{5pt}
\setlength{\emergencystretch}{3em}
\begin{document}
\section*{Drawing disconnected graphs on surfaces}\label{3321}
{\small UnsolvedMath ID 3321 · OPG-322 · Graph Theory · OpenGarden}
\subsection{Definitions and mathematical statement}
Let $G_1,G_2$ be finite graphs with disjoint vertex sets and no edge between them, and let $G=G_1\mathbin{\dot\cup}G_2$. Let $\Sigma$ be a connected closed surface. A drawing places vertices at distinct points and represents each edge by a simple arc between its endpoints. Edge interiors avoid vertices. Crossings are transverse intersections of edge interiors; there are finitely many, no three edge interiors meet at one point, and no tangencies or overlapping arc segments are allowed.

Use good drawings: adjacent edges do not cross and two independent edges cross at most once. Count each crossing point once. The surface crossing number $\operatorname{cr}_\Sigma(G)$ is the minimum number of crossings among such drawings, and a drawing attaining it is optimal. The number includes crossings between different components as well as crossings within one component.

The conjecture asks whether every optimal drawing of $G$ on $\Sigma$ has disjoint images of $G_1$ and $G_2$. Equivalently, no edge of the first graph should cross an edge of the second in any drawing attaining $\operatorname{cr}_\Sigma(G)$.

The word every is part of the question: constructing one optimum without intercomponent crossings would be a weaker result. No surface region or handle is assigned to a component in advance, and neither $G_1$ nor $G_2$ need be planar or connected. Disjointness concerns their entire drawn images, while crossings among edges within an individual component remain permissible and may be unavoidable.
\subsection{Short English statement}
In every drawing with the minimum possible crossings on a surface, must two disjoint graph components have disjoint images?
\subsection{Sources}
\begin{itemize}
\item[S1] ULAM.AI and the UnsolvedMath Contributors, supplied problems.json, record 3321 (OPG-322), OpenGarden. Catalogue identity and source transcription; CC BY 4.0.
\url{https://huggingface.co/datasets/ulamai/UnsolvedMath}.
\item[S2] Open Problem Garden, Drawing disconnected graphs on surfaces. Original problem and discussion.
\url{http://www.openproblemgarden.org/op/drawing_disconnected_graphs_on_surfaces}.
\end{itemize}
\end{document}
